% Chapter 8 — weave the graph: one build, every representation (kym_kg).

\gmenter{graph}
\note{Chapter 8: everything we learned becomes one graph. Four stops: Doge's corner of it, the
two languages it is published in, the shape of the whole graph, and how a build is published
without ever showing a half-built graph.}

\gmgo{graph-doge}{Doge, as dots and arrows}
\begin{frame}{Doge, step 6: all of it, as dots and arrows}
  \centering
  \resizebox{\linewidth}{!}{% generated by scripts/figures.py — do not edit
\begin{tikzpicture}[x=1cm, y=1cm,
 dot/.style={circle, inner sep=0pt, minimum size=#1},
 lab/.style={font=\fontsize{7.4}{8.4}\selectfont, text=gmInk, inner sep=.8pt},
 edge/.style={draw=gmBase, line width=.35pt}]
\node[dot=8mm, fill=kFrame] (doge) at (0.00,0.00) {};
\node[font=\scriptsize\bfseries, text=white] at (doge) {Doge};
\node[dot=3.2mm, fill=kFrame] (anc0) at (0.00,1.00) {};
\node[lab, anchor=west] at (anc0.east) {Interior Monologue Captioning};
\node[dot=3.2mm, fill=kFrame] (anc1) at (0.00,1.62) {};
\node[lab, anchor=west] at (anc1.east) {Image Macros};
\node[dot=3.2mm, fill=kFrame] (anc2) at (0.00,2.24) {};
\node[lab, anchor=west] at (anc2.east) {Memes};
\node[dot=2.8mm, fill=kFrame] (ch0) at (-1.70,2.60) {};
\node[lab, anchor=east] at (ch0.west) {Doge 2 / Caesar};
\node[dot=2.8mm, fill=kFrame] (ch1) at (-2.00,2.17) {};
\node[lab, anchor=east] at (ch1.west) {Ironic Doge Memes};
\node[dot=2.8mm, fill=kFrame] (ch2) at (-1.70,1.73) {};
\node[lab, anchor=east] at (ch2.west) {Shiba Inus / Shibes};
\node[dot=2.8mm, fill=kFrame] (ch3) at (-2.00,1.30) {};
\node[lab, anchor=east] at (ch3.west) {The Death Of Kabosu / Doge};
\node[dot=2.3mm, fill=kFrame] (sb0) at (-5.10,1.00) {};
\node[lab, anchor=east] at (sb0.west) {Oh You So Crazy};
\node[dot=2.3mm, fill=kFrame] (sb1) at (-5.40,0.68) {};
\node[lab, anchor=east] at (sb1.west) {Stop Das Gay};
\node[dot=2.3mm, fill=kFrame] (sb2) at (-5.10,0.36) {};
\node[lab, anchor=east] at (sb2.west) {Animals Talking In All Caps};
\node[dot=2.3mm, fill=kFrame] (sb3) at (-5.40,0.04) {};
\node[lab, anchor=east] at (sb3.west) {Laughing Lizard / hhhehehe};
\node[dot=2.3mm, fill=kFrame] (sb4) at (-5.10,-0.29) {};
\node[lab, anchor=east] at (sb4.west) {Captioned Stock Photos};
\node[dot=2.3mm, fill=kFrame] (sb5) at (-5.40,-0.61) {};
\node[lab, anchor=east] at (sb5.west) {Master Trole Kid};
\node[dot=2.3mm, fill=kFrame] (sb6) at (-5.10,-0.93) {};
\node[lab, anchor=east] at (sb6.west) {Snek};
\node[dot=2.3mm, fill=kFrame] (sb7) at (-5.40,-1.25) {};
\node[lab, anchor=east] at (sb7.west) {Stop It Son, You Are Doing Me A Frighten};
\node[dot=2.2mm, fill=kWikidata] (wdQ384060) at (3.80,2.55) {};
\node[lab, anchor=west] at (wdQ384060.east) {Tumblr};
\node[dot=2.2mm, fill=kWikidata] (wdQ39315) at (3.80,1.99) {};
\node[lab, anchor=west] at (wdQ39315.east) {Shiba Inu};
\node[dot=2.2mm, fill=kWikidata] (wdQ23486479) at (3.80,1.43) {};
\node[lab, anchor=west] at (wdQ23486479.east) {Kabosu};
\node[dot=2.2mm, fill=kWikidata] (wdQ15613810) at (3.80,0.87) {};
\node[lab, anchor=west] at (wdQ15613810.east) {Doge};
\node[dot=2.2mm, fill=kWikidata] (wdQ261197) at (3.80,0.31) {};
\node[lab, anchor=west] at (wdQ261197.east) {monologue};
\node[dot=2.2mm, fill=kWikidata] (wdQ868492) at (3.80,-0.25) {};
\node[lab, anchor=west] at (wdQ868492.east) {Comic Sans MS};
\node[dot=2.2mm, fill=kWikidata] (wdQ144) at (5.85,2.55) {};
\node[lab, anchor=west] at (wdQ144.east) {dog};
\node[dot=2.2mm, fill=kWikidata] (wdQ131723) at (5.85,1.99) {};
\node[lab, anchor=west] at (wdQ131723.east) {Bitcoin};
\node[dot=2.2mm, fill=kWikidata] (wdQ1136) at (5.85,1.43) {};
\node[lab, anchor=west] at (wdQ1136.east) {Reddit};
\node[dot=2.2mm, fill=kWikidata] (wdQ15377916) at (5.85,0.87) {};
\node[lab, anchor=west] at (wdQ15377916.east) {Dogecoin};
\node[dot=2.2mm, fill=kWikidata] (wdQ5287) at (5.85,0.31) {};
\node[lab, anchor=west] at (wdQ5287.east) {Japanese};
\node[dot=2.7mm, fill=kTemplate] (t0) at (2.45,-0.75) {};
\node[lab, anchor=west] at (t0.east) {Doge};
\node[dot=2.7mm, fill=kTemplate] (t1) at (2.45,-1.06) {};
\node[lab, anchor=west] at (t1.east) {Average Fan vs. Average Enjoyer};
\node[dot=2.7mm, fill=kTemplate] (t2) at (2.45,-1.38) {};
\node[lab, anchor=west] at (t2.east) {Doge 2};
\node[dot=2.7mm, fill=kTemplate] (t3) at (2.45,-1.69) {};
\node[lab, anchor=west] at (t3.east) {doge then and now};
\node[dot=2.7mm, fill=kTemplate] (t4) at (2.45,-2.01) {};
\node[lab, anchor=west] at (t4.east) {Dogecoin};
\node[dot=2.7mm, fill=kTemplate] (t5) at (2.45,-2.32) {};
\node[lab, anchor=west] at (t5.east) {Swole Doge vs Cheems};
\node[dot=2.7mm, fill=kTemplate] (t6) at (2.45,-2.64) {};
\node[lab, anchor=west] at (t6.east) {Buff Doge vs. Cheems};
\node[dot=2.7mm, fill=kTemplate] (t7) at (2.45,-2.95) {};
\node[lab, anchor=west] at (t7.east) {Multi Doge};
\node[dot=2.2mm, fill=kWikidata] (wdQ33767) at (6.55,-1.10) {};
\node[lab, anchor=west] at (wdQ33767.east) {hand};
\node[dot=2.2mm, fill=kWikidata] (wdQ131514) at (6.55,-1.63) {};
\node[lab, anchor=west] at (wdQ131514.east) {couch};
\node[dot=2.2mm, fill=kWikidata] (wdQ317521) at (6.55,-2.17) {};
\node[lab, anchor=west] at (wdQ317521.east) {Elon Musk};
\node[dot=2.2mm, fill=kWikidata] (wdQ110983191) at (6.55,-2.70) {};
\node[lab, anchor=west] at (wdQ110983191.east) {Balltze};
\node[dot=2.5mm, fill=kConcept] (ty0) at (-1.60,-1.55) {};
\node[lab, anchor=east] at (ty0.west) {slang};
\node[dot=2.5mm, fill=kConcept] (ty1) at (-1.60,-1.81) {};
\node[lab, anchor=east] at (ty1.west) {character};
\node[dot=2.5mm, fill=kConcept] (ty2) at (-1.60,-2.08) {};
\node[lab, anchor=east] at (ty2.west) {animal};
\node[dot=2.5mm, fill=kConcept] (ty3) at (-1.60,-2.34) {};
\node[lab, anchor=east] at (ty3.west) {image-macro};
\node[dot=2.5mm, fill=kConcept] (ty4) at (-1.60,-2.60) {};
\node[lab, anchor=east] at (ty4.west) {exploitable};
\node[dot=1.6mm, fill=kConcept] (tg0) at (-3.60,-1.55) {};
\node[dot=1.6mm, fill=kConcept] (tg1) at (-3.60,-1.91) {};
\node[dot=1.6mm, fill=kConcept] (tg2) at (-3.60,-2.27) {};
\node[dot=1.6mm, fill=kConcept] (tg3) at (-3.60,-2.63) {};
\node[dot=1.6mm, fill=kConcept] (tg4) at (-4.02,-1.55) {};
\node[dot=1.6mm, fill=kConcept] (tg5) at (-4.02,-1.91) {};
\node[dot=1.6mm, fill=kConcept] (tg6) at (-4.02,-2.27) {};
\node[dot=1.6mm, fill=kConcept] (tg7) at (-4.02,-2.63) {};
\node[dot=1.6mm, fill=kConcept] (tg8) at (-4.44,-1.55) {};
\node[dot=1.6mm, fill=kConcept] (tg9) at (-4.44,-1.91) {};
\node[dot=1.6mm, fill=kConcept] (tg10) at (-4.44,-2.27) {};
\node[dot=1.6mm, fill=kConcept] (tg11) at (-4.44,-2.63) {};
\node[dot=1.6mm, fill=kConcept] (tg12) at (-4.86,-1.55) {};
\node[dot=1.6mm, fill=kConcept] (tg13) at (-4.86,-1.91) {};
\node[dot=1.6mm, fill=kConcept] (tg14) at (-4.86,-2.27) {};
\node[dot=1.6mm, fill=kConcept] (tg15) at (-4.86,-2.63) {};
\node[dot=1.6mm, fill=kConcept] (tg16) at (-5.28,-1.55) {};
\node[dot=1.6mm, fill=kConcept] (tg17) at (-5.28,-1.91) {};
\node[dot=1.6mm, fill=kConcept] (tg18) at (-5.28,-2.27) {};
\node[dot=1.6mm, fill=kConcept] (tg19) at (-5.28,-2.63) {};
\node[dot=1.6mm, fill=kConcept] (tg20) at (-5.70,-1.55) {};
\node[lab, text=gmInk2, anchor=north] at (-4.65,-2.95) {21 tags};
\node[dot=2.0mm, fill=kEvent] (ev0) at (-1.05,-2.30) {};
\node[dot=2.0mm, fill=kEvent] (ev1) at (-0.91,-2.34) {};
\node[dot=2.0mm, fill=kEvent] (ev2) at (-0.77,-2.37) {};
\node[dot=2.0mm, fill=kEvent] (ev3) at (-0.63,-2.41) {};
\node[dot=2.0mm, fill=kEvent] (ev4) at (-0.49,-2.43) {};
\node[dot=2.0mm, fill=kEvent] (ev5) at (-0.35,-2.46) {};
\node[dot=2.0mm, fill=kEvent] (ev6) at (-0.21,-2.47) {};
\node[dot=2.0mm, fill=kEvent] (ev7) at (-0.07,-2.48) {};
\node[dot=2.0mm, fill=kEvent] (ev8) at (0.07,-2.48) {};
\node[dot=2.0mm, fill=kEvent] (ev9) at (0.21,-2.47) {};
\node[dot=2.0mm, fill=kEvent] (ev10) at (0.35,-2.46) {};
\node[dot=2.0mm, fill=kEvent] (ev11) at (0.49,-2.43) {};
\node[dot=2.0mm, fill=kEvent] (ev12) at (0.63,-2.41) {};
\node[dot=2.0mm, fill=kEvent] (ev13) at (0.77,-2.37) {};
\node[dot=2.0mm, fill=kEvent] (ev14) at (0.91,-2.34) {};
\node[dot=2.0mm, fill=kEvent] (ev15) at (1.05,-2.30) {};
\node[lab, text=gmInk2, anchor=north] at (0,-2.62) {16 events, one story};
\begin{pgfonlayer}{background}
\draw[edge, draw=kFrame!60!gmSurface, line width=.7pt, -{Stealth[length=1.5mm]}] (doge) -- (anc0);
\draw[edge, draw=kFrame!60!gmSurface, line width=.7pt, -{Stealth[length=1.5mm]}] (anc0) -- (anc1);
\draw[edge, draw=kFrame!60!gmSurface, line width=.7pt, -{Stealth[length=1.5mm]}] (anc1) -- (anc2);
\draw[edge, draw=kFrame!60!gmSurface, -{Stealth[length=1.4mm]}] (ch0) -- (doge);
\draw[edge, draw=kFrame!60!gmSurface, -{Stealth[length=1.4mm]}] (ch1) -- (doge);
\draw[edge, draw=kFrame!60!gmSurface, -{Stealth[length=1.4mm]}] (ch2) -- (doge);
\draw[edge, draw=kFrame!60!gmSurface, -{Stealth[length=1.4mm]}] (ch3) -- (doge);
\draw[edge, dashed, draw=kFrame!45!gmSurface] (doge) -- (sb0);
\draw[edge, dashed, draw=kFrame!45!gmSurface] (doge) -- (sb1);
\draw[edge, dashed, draw=kFrame!45!gmSurface] (doge) -- (sb2);
\draw[edge, dashed, draw=kFrame!45!gmSurface] (doge) -- (sb3);
\draw[edge, dashed, draw=kFrame!45!gmSurface] (doge) -- (sb4);
\draw[edge, dashed, draw=kFrame!45!gmSurface] (doge) -- (sb5);
\draw[edge, dashed, draw=kFrame!45!gmSurface] (doge) -- (sb6);
\draw[edge, dashed, draw=kFrame!45!gmSurface] (doge) -- (sb7);
\draw[edge] (doge) -- (wdQ384060);
\draw[edge] (doge) -- (wdQ39315);
\draw[edge] (doge) -- (wdQ23486479);
\draw[edge] (doge) -- (wdQ15613810);
\draw[edge] (doge) -- (wdQ261197);
\draw[edge] (doge) -- (wdQ868492);
\draw[edge] (doge) -- (wdQ144);
\draw[edge] (doge) -- (wdQ131723);
\draw[edge] (doge) -- (wdQ1136);
\draw[edge] (doge) -- (wdQ15377916);
\draw[edge] (doge) -- (wdQ5287);
\draw[edge] (doge) -- (t0);
\draw[edge] (doge) -- (t1);
\draw[edge] (doge) -- (t2);
\draw[edge] (doge) -- (t3);
\draw[edge] (doge) -- (t4);
\draw[edge] (doge) -- (t5);
\draw[edge] (doge) -- (t6);
\draw[edge] (doge) -- (t7);
\draw[edge, draw=kTemplate!55!gmSurface] (t2) -- (wdQ33767);
\draw[edge, draw=kTemplate!55!gmSurface] (t2) -- (wdQ131514);
\draw[edge, draw=kTemplate!55!gmSurface] (t4) -- (wdQ317521);
\draw[edge, draw=kTemplate!55!gmSurface] (t5) -- (wdQ110983191);
\draw[edge, draw=kTemplate!55!gmSurface] (t6) -- (wdQ110983191);
\draw[edge, draw=kConcept!40!gmSurface] (doge) -- (ty0);
\draw[edge, draw=kConcept!40!gmSurface] (doge) -- (ty1);
\draw[edge, draw=kConcept!40!gmSurface] (doge) -- (ty2);
\draw[edge, draw=kConcept!40!gmSurface] (doge) -- (ty3);
\draw[edge, draw=kConcept!40!gmSurface] (doge) -- (ty4);
\draw[edge, draw=kConcept!25!gmSurface, line width=.25pt] (doge) -- (tg0);
\draw[edge, draw=kConcept!25!gmSurface, line width=.25pt] (doge) -- (tg1);
\draw[edge, draw=kConcept!25!gmSurface, line width=.25pt] (doge) -- (tg2);
\draw[edge, draw=kConcept!25!gmSurface, line width=.25pt] (doge) -- (tg3);
\draw[edge, draw=kConcept!25!gmSurface, line width=.25pt] (doge) -- (tg4);
\draw[edge, draw=kConcept!25!gmSurface, line width=.25pt] (doge) -- (tg5);
\draw[edge, draw=kConcept!25!gmSurface, line width=.25pt] (doge) -- (tg6);
\draw[edge, draw=kConcept!25!gmSurface, line width=.25pt] (doge) -- (tg7);
\draw[edge, draw=kConcept!25!gmSurface, line width=.25pt] (doge) -- (tg8);
\draw[edge, draw=kConcept!25!gmSurface, line width=.25pt] (doge) -- (tg9);
\draw[edge, draw=kConcept!25!gmSurface, line width=.25pt] (doge) -- (tg10);
\draw[edge, draw=kConcept!25!gmSurface, line width=.25pt] (doge) -- (tg11);
\draw[edge, draw=kConcept!25!gmSurface, line width=.25pt] (doge) -- (tg12);
\draw[edge, draw=kConcept!25!gmSurface, line width=.25pt] (doge) -- (tg13);
\draw[edge, draw=kConcept!25!gmSurface, line width=.25pt] (doge) -- (tg14);
\draw[edge, draw=kConcept!25!gmSurface, line width=.25pt] (doge) -- (tg15);
\draw[edge, draw=kConcept!25!gmSurface, line width=.25pt] (doge) -- (tg16);
\draw[edge, draw=kConcept!25!gmSurface, line width=.25pt] (doge) -- (tg17);
\draw[edge, draw=kConcept!25!gmSurface, line width=.25pt] (doge) -- (tg18);
\draw[edge, draw=kConcept!25!gmSurface, line width=.25pt] (doge) -- (tg19);
\draw[edge, draw=kConcept!25!gmSurface, line width=.25pt] (doge) -- (tg20);
\draw[edge, draw=kEvent!40!gmSurface, line width=.3pt] (doge) -- (ev0);
\draw[edge, draw=kEvent!40!gmSurface, line width=.3pt] (doge) -- (ev1);
\draw[draw=kEvent, line width=.5pt, -{Stealth[length=1mm]}] (ev0) -- (ev1);
\draw[edge, draw=kEvent!40!gmSurface, line width=.3pt] (doge) -- (ev2);
\draw[draw=kEvent, line width=.5pt, -{Stealth[length=1mm]}] (ev1) -- (ev2);
\draw[edge, draw=kEvent!40!gmSurface, line width=.3pt] (doge) -- (ev3);
\draw[draw=kEvent, line width=.5pt, -{Stealth[length=1mm]}] (ev2) -- (ev3);
\draw[edge, draw=kEvent!40!gmSurface, line width=.3pt] (doge) -- (ev4);
\draw[draw=kEvent, line width=.5pt, -{Stealth[length=1mm]}] (ev3) -- (ev4);
\draw[edge, draw=kEvent!40!gmSurface, line width=.3pt] (doge) -- (ev5);
\draw[draw=kEvent, line width=.5pt, -{Stealth[length=1mm]}] (ev4) -- (ev5);
\draw[edge, draw=kEvent!40!gmSurface, line width=.3pt] (doge) -- (ev6);
\draw[draw=kEvent, line width=.5pt, -{Stealth[length=1mm]}] (ev5) -- (ev6);
\draw[edge, draw=kEvent!40!gmSurface, line width=.3pt] (doge) -- (ev7);
\draw[draw=kEvent, line width=.5pt, -{Stealth[length=1mm]}] (ev6) -- (ev7);
\draw[edge, draw=kEvent!40!gmSurface, line width=.3pt] (doge) -- (ev8);
\draw[draw=kEvent, line width=.5pt, -{Stealth[length=1mm]}] (ev7) -- (ev8);
\draw[edge, draw=kEvent!40!gmSurface, line width=.3pt] (doge) -- (ev9);
\draw[draw=kEvent, line width=.5pt, -{Stealth[length=1mm]}] (ev8) -- (ev9);
\draw[edge, draw=kEvent!40!gmSurface, line width=.3pt] (doge) -- (ev10);
\draw[draw=kEvent, line width=.5pt, -{Stealth[length=1mm]}] (ev9) -- (ev10);
\draw[edge, draw=kEvent!40!gmSurface, line width=.3pt] (doge) -- (ev11);
\draw[draw=kEvent, line width=.5pt, -{Stealth[length=1mm]}] (ev10) -- (ev11);
\draw[edge, draw=kEvent!40!gmSurface, line width=.3pt] (doge) -- (ev12);
\draw[draw=kEvent, line width=.5pt, -{Stealth[length=1mm]}] (ev11) -- (ev12);
\draw[edge, draw=kEvent!40!gmSurface, line width=.3pt] (doge) -- (ev13);
\draw[draw=kEvent, line width=.5pt, -{Stealth[length=1mm]}] (ev12) -- (ev13);
\draw[edge, draw=kEvent!40!gmSurface, line width=.3pt] (doge) -- (ev14);
\draw[draw=kEvent, line width=.5pt, -{Stealth[length=1mm]}] (ev13) -- (ev14);
\draw[edge, draw=kEvent!40!gmSurface, line width=.3pt] (doge) -- (ev15);
\draw[draw=kEvent, line width=.5pt, -{Stealth[length=1mm]}] (ev14) -- (ev15);
\end{pgfonlayer}
\end{tikzpicture}
}\par\vskip1mm
  % generated by scripts/figures.py — do not edit
\begin{tikzpicture}[x=1cm]
\fill[kFrame] (0.00,0) circle (.75mm); \node[anchor=west, font=\scriptsize, text=gmInk2] at (0.10,0) {frames};
\fill[kWikidata] (2.10,0) circle (.75mm); \node[anchor=west, font=\scriptsize, text=gmInk2] at (2.20,0) {Wikidata items};
\fill[kTemplate] (4.20,0) circle (.75mm); \node[anchor=west, font=\scriptsize, text=gmInk2] at (4.30,0) {templates};
\fill[kEvent] (6.30,0) circle (.75mm); \node[anchor=west, font=\scriptsize, text=gmInk2] at (6.40,0) {events};
\fill[kConcept] (8.40,0) circle (.75mm); \node[anchor=west, font=\scriptsize, text=gmInk2] at (8.50,0) {types and tags};
\end{tikzpicture}

  \note{
  \begin{itemize}
    \item Layman version: every fact from the previous chapters is now an arrow from Doge.
    \item Up: the series it belongs to --- Interior Monologue Captioning, Image Macros, Memes. Upper
      left: the memes that grew out of it (Ironic Doge Memes, Shiba Inus, Doge 2, the death of
      Kabosu). Left, dashed: its siblings in the same series.
    \item Right: the Wikidata items its text names. Lower right: its templates, and what their
      images show --- some are the same items the text names (Shiba Inu, Doge): text and picture
      agree, and the graph shows it.
    \item Bottom: its 16 events, one chain. Lower left: its types and 22 tags.
    \item Drawn from the live graph, version \NVersion\ (published \NPublished), which holds Doge once,
      at his new address (chapter 2).
  \end{itemize}}
\end{frame}

\gmgo{graph-languages}{Two languages}
\begin{frame}{The same facts, in two languages}
  \begin{columns}[t]
    \column{.55\textwidth}
      {\large\bfseries\color{gmInk}RDF}\enspace{\normalsize\color{gmInk2}asked with SPARQL}\par\vskip3mm
      {\normalsize\ttfamily\linespread{1.25}\selectfont
      kym:doge \textcolor{gmInk2}{a} m4s:MediaFrame, kym:Meme ;\\
      \hspace*{4mm}m4s:year 2010 ;\\
      \hspace*{4mm}m4s:fromTags wd:Q144 ;\\
      \hspace*{4mm}mk:hasTemplate mk:template/8072285 .\\[3mm]
      \textcolor{gmInk2}{<<} kym:doge m4s:fromTags wd:Q144 \textcolor{gmInk2}{>>}\\
      \hspace*{4mm}mk:mentionText "dog" ;\\
      \hspace*{4mm}mk:linkScore 0.616 .}
    \column{.42\textwidth}
      {\large\bfseries\color{gmInk}A property graph}\enspace{\normalsize\color{gmInk2}asked with Cypher}\par\vskip3mm
      {\normalsize\ttfamily\linespread{1.25}\selectfont
      MATCH (f:Frame \{label: 'Doge'\})\\
      \hspace*{4mm}-[r:fromTags]->\\
      \hspace*{4mm}(e \{qid: 'Q144'\})\\
      RETURN f.year,\\
      \hspace*{4mm}r.mention\_texts,\\
      \hspace*{4mm}r.link\_scores}
  \end{columns}
  \vskip5mm
  {\large\color{gmInk2}One build, two languages: \textbf{\color{gmInk}IMKG's words} in RDF; the details \textbf{\color{gmInk}on the arrow} in both.}
  \note{
  \begin{itemize}
    \item The graph is published twice, from the same data: as RDF (the W3C standard, queried with
      SPARQL, in Fuseki) and as a property graph (Neo4j, queried with Cypher).
    \item Left: Doge in RDF. m4s: and kym: are IMKG's own vocabulary; a series is skos:broader ---
      exactly what IMKG emits, so IMKG's queries still run. mk: is MemeAtlas's namespace, for
      what IMKG has no word for (templates, events, curation \dots). Doge is written the way IMKG
      names him, kym:doge; since 7.1.0 his IRI is KYM's new address, \dots/sensitive/memes/doge
      (chapter 2) --- the triples are his, word for word.
    \item The double angle brackets are RDF-star: a statement about a statement. ``Doge's tags
      mention dog'' carries the words it was read from and the link's score.
    \item Right: the same fact in Cypher; a property-graph edge carries those details as its own
      properties.
  \end{itemize}}
\end{frame}

\gmgo{graph-meta}{The shape of the graph}
\begin{frame}{The graph of the graph}
  \begin{columns}[c]
    \column{.8\textwidth}\centering
      \resizebox{\linewidth}{!}{% generated by scripts/figures.py — do not edit
\begin{tikzpicture}[x=1cm, y=1cm, lab/.style={font=\fontsize{8}{9}\selectfont, text=gmInk2, fill=gmSurface, inner sep=.7pt}]
\node[circle, fill=kFrame, minimum size=0.97cm, inner sep=0pt] (k-frame) at (4.3,3.05) {};
\node[font=\fontsize{8.5}{9.5}\selectfont, text=gmInk, align=center, below=.5mm of k-frame, fill=gmSurface, inner sep=.6pt] {\textbf{frame} 24,291\\{\color{gmInk2}$\circlearrowleft$ siblings, links 856k}};
\node[circle, fill=kEvent, minimum size=1.77cm, inner sep=0pt] (k-event) at (8.5,3.05) {};
\node[font=\fontsize{8.5}{9.5}\selectfont, text=gmInk, align=center, below=.5mm of k-event, fill=gmSurface, inner sep=.6pt] {\textbf{event} 142,787\\{\color{gmInk2}$\circlearrowleft$ nextInStory 134k}};
\node[circle, fill=kDoc, minimum size=2.24cm, inner sep=0pt] (k-external_ref) at (11.7,5.5) {};
\node[font=\fontsize{8.5}{9.5}\selectfont, text=gmInk, align=center, below=.5mm of k-external_ref, fill=gmSurface, inner sep=.6pt] {\textbf{external page} 257,031};
\node[circle, fill=kDoc, minimum size=2.03cm, inner sep=0pt] (k-image) at (11.7,0.6) {};
\node[font=\fontsize{8.5}{9.5}\selectfont, text=gmInk, align=center, below=.5mm of k-image, fill=gmSurface, inner sep=.6pt] {\textbf{image} 202,390};
\node[circle, fill=kWikidata, minimum size=1.00cm, inner sep=0pt] (k-wikidata_entity) at (0.9,5.75) {};
\node[font=\fontsize{8.5}{9.5}\selectfont, text=gmInk, align=center, below=.5mm of k-wikidata_entity, fill=gmSurface, inner sep=.6pt] {\textbf{Wikidata item} 26,897};
\node[circle, fill=kTemplate, minimum size=0.99cm, inner sep=0pt] (k-template) at (6.6,0.0) {};
\node[font=\fontsize{8.5}{9.5}\selectfont, text=gmInk, align=center, below=.5mm of k-template, fill=gmSurface, inner sep=.6pt] {\textbf{template} 26,868};
\node[circle, fill=kStub, minimum size=0.72cm, inner sep=0pt] (k-frame_stub) at (6.0,6.05) {};
\node[font=\fontsize{8.5}{9.5}\selectfont, text=gmInk, align=center, below=.5mm of k-frame_stub, fill=gmSurface, inner sep=.6pt] {\textbf{frame stub} 7,631};
\node[circle, fill=kConcept, minimum size=1.56cm, inner sep=0pt] (k-tag_concept) at (0.7,2.75) {};
\node[font=\fontsize{8.5}{9.5}\selectfont, text=gmInk, align=center, below=.5mm of k-tag_concept, fill=gmSurface, inner sep=.6pt] {\textbf{tag} 101,883\\{\color{gmInk2}$\circlearrowleft$ coOccursWith 5k}};
\node[circle, fill=kConcept, minimum size=0.44cm, inner sep=0pt] (k-entry_type_concept) at (1.3,0.35) {};
\node[font=\fontsize{8.5}{9.5}\selectfont, text=gmInk, align=center, below=.5mm of k-entry_type_concept, fill=gmSurface, inner sep=.6pt] {\textbf{entry type} 119};
\node[circle, fill=kConcept, minimum size=0.67cm, inner sep=0pt] (k-origin_concept) at (3.4,-0.45) {};
\node[font=\fontsize{8.5}{9.5}\selectfont, text=gmInk, align=center, below=.5mm of k-origin_concept, fill=gmSurface, inner sep=.6pt] {\textbf{origin} 5,703};
\node[circle, fill=kConcept, minimum size=0.44cm, inner sep=0pt] (k-region_concept) at (3.2,6.05) {};
\node[font=\fontsize{8.5}{9.5}\selectfont, text=gmInk, align=center, below=.5mm of k-region_concept, fill=gmSurface, inner sep=.6pt] {\textbf{region} 110};
\begin{pgfonlayer}{background}
\draw[gmBase, line width=2.56pt, -{Stealth[length=2.9mm]}] (k-frame) -- node[lab, sloped, above, pos=0.38] {citesExternal 244k} (k-external_ref);
\draw[gmBase, line width=2.54pt, -{Stealth[length=2.9mm]}] (k-frame) -- node[lab, sloped, above, pos=0.5] {hasTag 226k} (k-tag_concept);
\draw[gmBase, line width=2.50pt, -{Stealth[length=2.9mm]}] (k-frame) to[bend right=10] node[lab, sloped, above, pos=0.3] {hasImage 193k} (k-image);
\draw[gmBase, line width=2.49pt, -{Stealth[length=2.8mm]}] (k-frame) -- node[lab, sloped, above, pos=0.4] {about, tags, title 181k} (k-wikidata_entity);
\draw[gmBase, line width=2.42pt, -{Stealth[length=2.8mm]}] (k-frame) to[bend left=14] node[lab, sloped, above, pos=0.5] {hasEvent 143k} (k-event);
\draw[gmBase, line width=2.41pt, -{Stealth[length=2.8mm]}] (k-event) -- node[lab, sloped, above, pos=0.38] {citations, embeds 138k} (k-external_ref);
\draw[gmBase, line width=2.34pt, -{Stealth[length=2.8mm]}] (k-event) -- node[lab, sloped, above, pos=0.45] {eventImage 104k} (k-image);
\draw[gmBase, line width=2.23pt, -{Stealth[length=2.7mm]}] (k-event) to[bend left=14] node[lab, sloped, above, pos=0.5] {eventLink 67k} (k-frame);
\draw[gmBase, line width=2.19pt, -{Stealth[length=2.7mm]}] (k-frame) -- node[lab, sloped, above, pos=0.45] {citesMediaFrame 58k} (k-frame_stub);
\draw[gmBase, line width=2.12pt, -{Stealth[length=2.7mm]}] (k-template) to[bend left=26] node[lab, sloped, above, pos=0.7] {fromImage 45k} (k-wikidata_entity);
\draw[gmBase, line width=2.05pt, -{Stealth[length=2.6mm]}] (k-frame) -- node[lab, sloped, above, pos=0.55] {hasEntryType 34k} (k-entry_type_concept);
\draw[gmBase, line width=2.02pt, -{Stealth[length=2.6mm]}] (k-frame) -- node[lab, sloped, above, pos=0.5] {hasTemplate 30k} (k-template);
\draw[gmBase, line width=1.99pt, -{Stealth[length=2.6mm]}] (k-template) to[bend right=24] node[lab, sloped, above, pos=0.55] {imgflipPage 27k} (k-external_ref);
\draw[gmBase, line width=1.97pt, -{Stealth[length=2.6mm]}] (k-template) -- node[lab, sloped, above, pos=0.5] {templateImage 25k} (k-image);
\draw[gmBase, line width=1.96pt, -{Stealth[length=2.6mm]}] (k-frame) -- (k-origin_concept);
\draw[gmBase, line width=1.94pt, -{Stealth[length=2.6mm]}] (k-event) to[bend right=12] (k-frame_stub);
\draw[gmBase, line width=1.58pt, -{Stealth[length=2.4mm]}] (k-frame) -- (k-region_concept);
\end{pgfonlayer}
\end{tikzpicture}
}
    \column{.18\textwidth}
      \gmbig{\NNodesShort}{nodes}\par\vskip3mm
      \gmbig{\NEdgesShort}{edges}\par\vskip3mm
      \gmbig{\NTriplesShort}{RDF triples}\par\vskip3mm
      {\normalsize\color{gmInk2}version \NVersion,\\\NRelTypes\ kinds of arrow}
  \end{columns}
  \note{
  \begin{itemize}
    \item Every kind of node, every kind of edge between two kinds, sized by how many there are
      (area: nodes; width: edges, on a log scale) --- the schema of the graph, measured, not drawn
      on a whiteboard.
    \item The frame is the hub. The biggest flows: Wikidata's statements between items (696k, new
      in 7.1.0, the loop under ``Wikidata item''), siblings (631k, derived from the series), links
      to external pages (236k), tags (218k), images (185k), events (137k) and the story chain
      between events (119k).
    \item The live graph, version \NVersion, published \NPublished: \NNodes\ nodes, \NEdges\ edges,
      \NTriplesShort\ RDF triples, \NRelTypes\ kinds of relation --- one of them Wikidata's statements, in
      \NStatementProps\ properties.
      Most Wikidata items are values of those statements (Kabosu \emph{is a} dog: ``dog'' is a node):
      of the \NNodesWikidata\ items, the disk counts the \NWikidataLinked\ that entries and templates
      link to.
  \end{itemize}}
\end{frame}

\gmgo{graph-stores}{One build, four stores}
\begin{frame}{One build, four stores, checked before publishing}
  \centering
  \begin{tikzpicture}[x=1cm, y=1cm,
      st/.style={rounded corners=2mm, draw=gm@col@graph, line width=1.2pt, fill=gmSurface, minimum width=2.75cm,
                 minimum height=1.35cm, align=center, font=\large\bfseries},
      ck/.style={rounded corners=2mm, fill=gm@col@graph, text=white, minimum width=3.2cm, minimum height=1.35cm,
                 align=center, font=\large\bfseries},
      go/.style={-{Stealth[length=3mm]}, line width=1.4pt, draw=gmBase, shorten <=1mm, shorten >=1mm}]
    \node[st, fill=gm@col@graph!8!gmSurface] (b) at (0,0) {build\\[-.3mm]{\normalsize\mdseries\color{gmInk2} every entry, in chunks}};
    \node[st] (m) at (3.55,0) {Mongo\\[-.3mm]{\normalsize\mdseries\color{gmInk2} the authority}};
    \node[st] (f) at (7.05,1.45) {files\\[-.3mm]{\normalsize\mdseries\color{gmInk2} graph.nt, CSV}};
    \node[st] (u) at (10.55,1.45) {Fuseki\\[-.3mm]{\normalsize\mdseries\color{gmInk2} RDF, SPARQL}};
    \node[st] (n) at (7.05,-1.45) {Neo4j\\[-.3mm]{\normalsize\mdseries\color{gmInk2} Cypher}};
    \draw[go] (b) -- (m); \draw[go] (m) -- (f); \draw[go] (f) -- (u); \draw[go] (m) -- (n);
    \draw[gm@col@graph!45!gmSurface, line width=1.2pt, dash pattern=on 4pt off 3pt, rounded corners=3mm]
      (1.9,-2.45) rectangle (12.15,2.45);
    \node[anchor=north west, font=\small, text=gmInk2] at (2.05,2.35) {beside the live build};
    \node[ck] (v) at (10.55,-1.1) {verify\\[-.3mm]{\normalsize\mdseries all four agree}};
    \draw[go, draw=gm@col@graph!50!gmSurface] (n.east) -- (v.west|-n.east);
    \draw[go, draw=gm@col@graph!50!gmSurface] (u) -- (v);
  \end{tikzpicture}
  \vskip2mm
  \begin{tikzpicture}[x=1cm, y=1cm]
    \node[rounded corners=2mm, fill=gm@col@graph, text=white, inner sep=3mm, text width=13.4cm, align=left] {
      {\large\bfseries publish}\enspace\normalsize files, then Fuseki, then Neo4j, \textbf{Mongo last}; every pointer
      read back. Weekly, the RDF is derived a second way: equal but for \NTriplesDiff\ provenance triples.};
  \end{tikzpicture}
  \note{
  \begin{itemize}
    \item kym\_kg: build.py is the only producer of nodes and edges. It streams every entry, with
      its events, links and templates, in chunks, into Mongo --- the authority --- under a new build
      id, beside the published build. Readers never see a half-built graph: nobody reads the new
      id yet.
    \item From Mongo, one stream writes the files (graph.nt, the CSVs the mapping reads, the
      manifest); Fuseki loads the RDF from graph.nt; Neo4j is loaded from Mongo.
    \item Verify: Mongo, the manifest, the RML rows, Fuseki and Neo4j must hold the same counts;
      if not, the build fails and nothing moves.
    \item Publish flips pointers: followers first, the authority last, each read back --- a crash
      leaves a follower ahead of the authority, never behind, and the next run reconciles it. A
      torn publish fails loudly. The stores keep the live build and the one before, for rollback.
    \item Independently, kym\_kg\_validate re-derives the RDF from the CSVs through the YARRRML
      mapping (morph-kgc) and diffs it with graph.nt: \NRmlTriples\ against \NTriples\ --- the
      \NTriplesDiff\ extra are the build's own provenance triples.
  \end{itemize}}
\end{frame}
